\section{Case Study Setup} \label{sec:casestudysetup}
To study the effectiveness of our approaches for detecting performance regressions under varying workloads, we perform case studies on two open-source systems\footnote{Our experimental setup, workloads, and results are shared online \url{https://github.com/senseconcordia/EMSE2020Data} as a replication package.}. In this section, we first present the subject systems. Then, we present the workloads applied to the systems, the experimental environment, and performance issues. Finally, we present the choices of machine learning models that are used in our approaches. 

\subsection{Subject systems}
We evaluate our approaches on two open-source systems, namely OpenMRS and Apache James. We also discuss the successful application of our approach in System X, a large-scale industrial system that provides government-regulation related reporting services (in Section~\ref{sec:deploymentinindustry}). Apache James is a Java-based mail server. OpenMRS is an open-source health care system that supports customized medical records and it is wildly used in developing countries. The two open-source subject systems and the industrial system are all studied in prior research~\citep{Yao:2018:LSL:3184407.3184416, DBLP:conf/icst/GaoJBL16} and cover different domains, which ensures that our findings are not limited to a specific domain. The details of the two open-source subject systems and the industrial system are shown in Table~\ref{tab:subjects}. 

Both two open-source subject systems and the industrial system are CPU-intensive. Besides, CPU is usually the main contributor to server costs~\citep{Greenberg:2008:CCR:1496091.1496103}. Performance regression in terms of CPU would result in the need for more CPU resources to provide the same quality of service, thereby significantly increasing the cost of system operations. 
%For example, for our industrial system (System X), CPU usage is the main concern in performance regression detection.
Therefore, in this study, we use CPU as the performance metric for detecting performance regressions. 
%\heng{Moved the CPU usage information to the end of 4.2 Subject workload design, as 1) we have not talked about the different versions here, 2) we just claimed that these systems are cpu-intensive. }
%(summarized in Table~\ref{tab:cpu_utilization}). \ian{added} We would like to note that when we observe a higher CPU usage, it may not correspond to a performance regression, while it may be just because the system has a higher workload. The influence from the workload to the CPU usage is one of the reason that we leverage performance modeling to detect performance regressions in the field.



\begin{table}[tbh]
  \centering
  \caption{Overview of the open-source subject systems and the industrial system.}
    \begin{threeparttable}
    \begin{tabular}{c|c|c|r}
    \hline
    Subjects & Versions & Domains & SLOC (K) \\
    \hline
    OpenMRS & 2.1.4 & Medical & 67 \\
    
    Apache James & 2.3.2, 3.0M1, 3.0M2 & Mail Server & 37 \\
    
    System X\tnote{*} & 10 releases in 2019 & Commercial & \textgreater 2,000 \\
    \hline
    \end{tabular}%
    \begin{tablenotes}
    \item[*] We share our experience of applying our approaches on System X in Section~\ref{sec:deploymentinindustry}.
    \end{tablenotes}
    \end{threeparttable}
  \label{tab:subjects}%
\end{table}%


% Table generated by Excel2LaTeX from sheet '___3'
\begin{table*}[htbp]
  \centering
  \tiny
  \caption{Summary of the test actions with increased appearances (the ones with $\bullet$) in different load drivers.}
    \begin{tabular}{|m{16em}<{\centering}|m{4em}<{\centering}|m{4em}<{\centering}|m{4em}<{\centering}|m{4em}<{\centering}|m{4em}<{\centering}|}
    \hline
    \multicolumn{6}{|c|}{OpenMRS} \\
    \hline
    \multicolumn{1}{|c|}{Test actions} & Load driver 1 & Load driver 2 & Load driver 3 & Load driver 4 & Load driver 5 \\
    \hline
    Creation of patients & $\bullet$   &       & $\bullet$   &       & $\bullet$ \\
    \hline
    Deletion of patients & $\bullet$   &       & $\bullet$   &       & $\bullet$ \\
    \hline
    Searching for patients & $\bullet$   & $\bullet$   &       &       &  \\
    \hline
    Editing patients &       &       &       & $\bullet$   & $\bullet$ \\
    \hline
    Searching for concepts &       &       & $\bullet$   & $\bullet$   & $\bullet$ \\
    \hline
    Searching for encounters &       & $\bullet$   & $\bullet$   &       &  \\
    \hline
    Searching for observations &       &       & $\bullet$   & $\bullet$   & $\bullet$ \\
    \hline
    Searching for types of encounters &       & $\bullet$   &       & $\bullet$   & $\bullet$ \\
    \hline
    \multicolumn{6}{c}{} \\
    \hline
    \multicolumn{6}{|c|}{Apache James} \\
    \hline
    \multicolumn{1}{|c|}{Test actions} & Load driver 1 & Load driver 2 & Load driver 3 & Load driver 4 & Load driver 5 \\
    \hline
    Sending mails with short messages and without attachments &       &       &       & $\bullet$   &  \\
    \hline
    Sending mails with long messages and without attachments & $\bullet$   & $\bullet$   & $\bullet$   &       & $\bullet$ \\
    \hline
    Sending mails with short messages and with small attachments &       &       & $\bullet$   & $\bullet$   &  \\
    \hline
    Sending mails with short messages and with large attachments &       & $\bullet$   &       &       &  \\
    \hline
    Sending mails with long messages and with small attachments &       &       & $\bullet$   &       &  \\
    \hline
    Sending mails with long messages and with large attachments & $\bullet$   & $\bullet$   &       &       &  \\
    \hline
    Retrieving entire mails &       &       &       &       & $\bullet$ \\
    \hline
    Retrieving only the header of mails & $\bullet$   &       &       & $\bullet$   & $\bullet$ \\
    \hline
    \end{tabular}%
  \label{tab:workload_difference}%
\end{table*}%



\subsection{Subject workload design}
\subsubsection{OpenMRS}
OpenMRS provides a web-based interface and RESTFul services. We used the default OpenMRS demo database in our performance tests. The demo database contains data for over 5K patients and 476K observations. OpenMRS contains four typical scenarios: adding, deleting, searching, and editing operations. We designed different performance tests that are composed of eight various test actions, including 1) creation of patients, 2) deletion of patients, 3) searching for patients, 4) editing patients, 5) searching for concepts, 6) searching for encounters, 7) searching for observations, and 8) searching for types of encounters.
%\ian{should be i.e., not e.g., right? and we need to number the actions}e.g., searches of patients, concepts, encounters, and observations, and creation, deletion and edit of patient information.
Those performance tests are different in the ratio between the actions. 
In order to simulate a more realistic workload in the field, we added random controllers and random order controllers in JMeter to vary the workload. Moreover, we also simulated the variety of the number of users and activities in the field by setting random gaps between the repetitions of each user’s activities, randomizing the order of the user activities, and setting a different number of maximum concurrent users for different workloads at different times.

Furthermore, we designed a total of five JMeter-based performance tests, each of which consists of a mixture of random workloads. % as described above. 
In particular, to drive different workloads, for each of the load drivers, we pick several different test actions and put them in an extra JMeter loop controller that iterates a random number of times to increase their appearances in the workload. Table~\ref{tab:workload_difference} shows the detailed choice of the looped actions of each load driver. 
To make sure that the two versions of the subject systems have different workloads, for one version of the system (e.g., the old version), we run four JMeter tests together to exercise the system. For the other version (e.g., the new version) of the system, we run one additional JMeter test, i.e, a total of five JMeter tests, to ensure that the two versions have undergone different workloads.
%\lizhi{Through the analysis of results, it is found that test action \emph{deletion of patients} has the largest portion of changes ranging from $2.7\%$ to $8.3\%$}

Despite the variation in the workloads, we keep our systems not saturated and running in normal conditions. The  CPU  saturation  can  be  qualified as the average CPU load increases to a very large value (e.g., 90\%) and remains for a long period of time ~\citep{10.1117/12.705735}. When the system is saturated, the services of the systems are not provided normally and the relationship between the performance of a software system (e.g.,  CPU usage) and the runtime activities may change. Besides, real software systems are usually not running under saturated conditions.  
The details of the workload, such as the CPU usage values, are shared in our replication package. 

%\lizhi{ref. R1.1 
For the OpenMRS system, we do not have any evidence showing performance regressions of specific versions. Therefore, we manually inject several performance regressions in the system. The injected performance regressions are shown in Table~\ref{tab:workloaddeisign}. 
The original version, i.e., v0 of OpenMRS does not have any injected or known performance regressions. 
We separately injected four performance regressions (heavier DB request, additional I/O, constant delay and additional calculation) on the basis of version v0. These four versions are called v1, v2, v3, and v4, respectively. The details of the injected performance regressions of OpenMRS are explained below:
\begin{itemize}
    \item[$\bullet$] \textbf{v1: Injected heavier DB request.}  We increased the number of DB requests in the code responsible for accessing the database in OpenMRS. This performance issue will increase the CPU usage of the MySQL server.
    \item[$\bullet$] \textbf{v2: Added additional I/O access.} Since accessing I/O storage devices (e.g., hard drives) are usually slower than accessing memory, we added redundant logging statements to the source code that is frequently executed. The execution of the logging statements will introduce performance regression.
    \item[$\bullet$] \textbf{v3: Created constant delay.} We added the delay bug in the frequently executed code, which creates a blocking performance issue.
    \item[$\bullet$] \textbf{v4: Injected additional calculation.} We added additional calculations to the source code of OpenMRS that are frequently executed during the performance test.
\end{itemize}
%}

We deployed the OpenMRS on two machines, each with a configuration of Intel Core i5-2400 CPU (3.10GHz), 8 GB memory, and 512GB SATA hard drive. One machine is deployed as the application server and the other machine is deployed as the MySQL database server. We ran JMeter using the RESTFul API of OpenMRS on five extra machines with the same specification to simulate user workloads on the client side for five hours. Each machine hosts one JMeter instance with one type of workload. We used \emph{Pidstat}~\citep{pidstat} to monitor the CPU usage that is used as the performance metrics of this study.
To minimize the noise from the system warm-up and cool-down periods, we filtered out the data from the first and last half hour of running each workload. Thus, we only kept four hours of data from each performance test.

\subsubsection{Apache James}
Apache James is an open-source enterprise mail server. It contains two main actions: sending and retrieving mails. Those two actions can be further divided into many smaller actions, e.g., sending mails with or without attachments and retrieving an entire mail or only the header of the mail. In total, we built eight actions for the performance test, including 1) sending mails with short messages and without attachments, 2) sending mails with long messages and without attachments, 3) sending mails with short messages and with small attachments, 4) sending mails with short messages and with large attachments, 5) sending mails with long messages and with small attachments, 6) sending mails with long messages and with large attachments, 7) retrieving entire mails, and 8) retrieving only the header of mails.
Similar to OpenMRS, we also created a total of five different workloads, where each workload consists of a mixture of eight test actions with different ratios. 
We conduct performance tests with different versions of Apache James with varying workloads (i.e., 4 tests for one version and 5 tests for the other version). We used JMeter to create performance tests that exercise Apache James.
Different from OpenMRS, in which the performance regressions are manually injected, we performance-tested on three different versions of systems with known real-world performance issues for Apache James. Apart from the last stable version 2.3.2, there are multiple Milestone Releases for version 3.0. After checking the release notes on the Apache James website, we picked 3.0M1 and 3.0M2 to use in our study. These three selected versions have many bug fixes and performance improvements~\citep{Apache-James}. The same subjects are studied in prior research on performance testing~\citep{DBLP:conf/icst/GaoJBL16}. Table~\ref{tab:workloaddeisign} summarizes the details of performance improvements of three selected versions. 

We deployed Apache James on a server machine with an Intel Core i7-8700K CPU (3.70GHz), 16 GB of memory, and a 3TB SATA hard drive. We ran JMeter on five extra machines with a configuration of Intel Core i5-2400 CPU (3.10GHz), 8 GB memory and 320GB SATA hard drive. These JMeter instances generate multiple workloads for five hours. Similar to OpenMRS, we use \emph{Pidstat}~\citep{pidstat} to monitor the CPU usage as the performance metrics of this study. We filtered out the data from the first and last half hour of the tests to minimize the system noise.

% \subsubsection{System X}
% System X is deployed in production and is used by real users worldwide. We retrieved the execution logs and the corresponding CPU usage on a daily basis. System X is deployed in an internal production environment. Due to the NDA, we cannot reveal additional details about the hardware environment and the usage scenarios of System X.



\begin{table}[tbh]
  \centering
  \caption{Performance regressions in the studied open-source systems}
    \begin{tabular}{c|c|p{3.3cm}}
    \hline
    System & Versions & \multicolumn{1}{c}{Performance regressions} \\
    \hline
          & v0    & \multicolumn{1}{c}{Original version} \\
\cline{2-3}          & v1    & \multicolumn{1}{c}{Injected heavier DB request} \\
\cline{2-3}    OpenMRS & v2    & \multicolumn{1}{c}{Added additional I/O access} \\
\cline{2-3}          & v3    & \multicolumn{1}{c}{Created constant delay} \\
\cline{2-3}          & v4    & \multicolumn{1}{c}{Injected additional calculation} \\
    \hline
          & 2.3.2 & \multicolumn{1}{c}{Stable release version} \\
          \cline{2-3}
          Apache James & 3.0M1 & \multicolumn{1}{c}{Improved ActiveMQ spool efficiency~\citep{Apache-James}} \\
          \cline{2-3}
          & 3.0M2& \multicolumn{1}{c}{Improved large attachment handling efficiency~\citep{Apache-James}} \\
    \hline
    \end{tabular}%
  \label{tab:workloaddeisign}%
\end{table}%

\subsubsection{Generated workloads}

Our generated workloads are different across different software versions.
By conducting chi-square tests of independence between the number of different actions and the versions, we find that in both OpenMRS and Apache James, the number of different actions and the versions are not statistically independent (p-value $\ll$ 0.05), which means that the distribution of the actions are statistically different across versions. For example, for the OpenMRS system, the total number of all actions ranges from 10,673 to 22,392 across versions. In particular, the number of the \emph{deletion of patients} action has a range from 293 to 1,861 across all versions. For another example from Apache James, the number of the \emph{retrieving only the header of mails} action has a range from 411 to 25,532 across all versions. In addition, we provide the detailed numbers of each action that are sent to each version in our replication package. 




Table~\ref{tab:cpu_utilization} summarizes the average CPU utilization of each version of the systems under our generated dynamic workloads. We would like to note that when we observe a higher CPU usage, it may not correspond to a performance regression; instead, it may be just due to the system having a higher workload. The influence of the dynamic workload on the CPU usage is one of the reasons that we leverage performance modeling to detect performance regressions in the field.
The detailed CPU utilization information during the entire test time period is reported in the replication package.

% Table generated by Excel2LaTeX from sheet '___2'
\begin{table}[tbh]
  \centering
  \footnotesize
  \caption{Summary of average CPU utilization of each version in our case studies.}
 
    \begin{tabular}{c|c|c|c|c|c}
\hline    \multicolumn{6}{c}{OpenMRS} \\
\hline
    v0 (4w) & v0 (5w) & v1    & v2    & v3    & \multicolumn{1}{c}{v4} \\
    \hline
     52.03 & 52.24 & 55.55 & 53.76 & 53.05 & 55.11 \\
    \hline
    \end{tabular}%
\vspace{1em}

    \begin{tabular}{c|c|c|c}
\hline  \multicolumn{4}{c}{ Apache James} \\
\hline
     3.0M2 (4w) & 3.0M2 (5w) & 3.0M1 & 2.3.2 \\
    \hline
     0.48  & 0.51  & 3.54  & 4.89 \\
    \hline
    \end{tabular}%

  \label{tab:cpu_utilization}%
\end{table}%

\subsection{Subject models}
Our approach presented in Section~\ref{sec:approach} is not designed strictly for any particular type of machine learning models. In fact, practitioners may choose their preferred models. In this study, we study the use of six different types of models, namely linear regression, random forest, XGBoost, convolutional neural network (CNN), recurrent neural network (RNN) and long short-term memory (LSTM). In particular, for linear regression, random forest, and XGBoost, the inputs of the models are vectors whose values are the number of appearances of each log event in a time period. For CNN, RNN, and LSTM, we sort the logs in each time period by their corresponding time stamp to create a sequence as the input of the neural networks. Our XGBoost model is fine-tuned using the GridSearchCV~\citep{girdsearcv} method. Our CNN, RNN, and LSTM have three, four and four layers, respectively, and they are trained with five, five and ten epochs, respectively. The number of layers and epochs are manually fine-tuned to avoid overfitting.
